\documentclass[english, parskip=full]{scrartcl}
\usepackage[english]{babel}  % Sprache auf Deutsch 
\usepackage[utf8]{inputenc}  % Kodierung der Datei
\usepackage[left=2cm, right=2cm, top=2cm, bottom=3cm, footskip=1.5cm]{geometry}
\input{myscrarticle}
\renewcommand{\r}{\text{r}}
\renewcommand{\t}{\text{t}}
\renewcommand{\a}{\text{a}}
\renewcommand{\o}{\text{o}}
\renewcommand{\F}{\text{F}}
\newcommand{\hc}{\text{h.c.}}
\newcommand{\B}{\text{B}}
\newcommand{\R}{\text{R}}
\renewcommand{\L}{\text{L}}
\newcommand{\TODO}[1]{\color{red!60!black} #1\color{black}}
%\newcommand{\QnA}[1]{\color{blue!80!black} #1\color{black}}
\newcommand\QnA[1]{\bgroup\color{blue} #1\egroup}

\newcommand*{\titel}{Notes on Luttinger Liquids}
\newcommand*{\autor}{Joachim Schwardt}

\hypersetup{pdfauthor={\autor}, pdftitle={\titel}}

%\titlehead{titlehead}
\subject{Quantum Physics in 1D}
\title{\titel}
\author{\autor}
\date{\today}

\begin{document}

%\begin{titlepage}
\maketitle  % Titel setzen
%\tableofcontents  % Inhaltsverzeichnis setzen
%\end{titlepage}

\section{Analytical tools}
\subsection{Gaussian integrals}
A very important formula is
\begin{align}
I &= \int_{\mathbb{C}^N} \exp\kla*{-\brc*{u}{A}{u} + \brb*{h}{u} + \brb*{u}{h}  } \, \frac{\mathrm{d}u_1\mathrm{d}u_1^\ast}{2\pi\i} \dots \frac{\mathrm{d}u_N\mathrm{d}u_N^\ast}{2\pi\i} = \frac{\powe{ \brc*{h}{A^{-1}}{h} }}{\det A}. \label{eqn:exp minus u ast A u plus h ast u plus u ast h}
\end{align}
Here $\Ket{v} \equiv (v_1,\dots,v_N)^T$ and $\Bra{v} \equiv (v_1^\ast, \dots, v_N^\ast)$.
The proof of this formula is relatively simple for the special case of diagonalizable matrices satisfying $A = UDU^\dagger$.
Then one may substitute $u = Uv$ to get
\begin{align}
I &= \int_{\mathbb{C}^N} \exp\kla*{-\brc*{v}{D}{v} + \brc*{h}{U}{v} + \brc{v}{U^\dagger}{h}  } \, \frac{\mathrm{d}v_1\mathrm{d}v_1^\ast}{2\pi\i} \dots \frac{\mathrm{d}v_N\mathrm{d}v_N^\ast}{2\pi\i}.
\end{align}
Now all the $v$'s are decoupled and the integral decomposes into
\begin{align}
I &= \prod_{j=1}^N \int_{\mathbb{R}^2} \powe{-(v_\r - \i v_\i) \lambda_j (v_\r + \i v_\i) + (U^\dagger h)_j (v_\r + \i v_\i) + (v_\r - \i v_\i) (U^\dagger h)_j^\ast } \frac{\mathrm{d}v_\r\mathrm{d}v_\i}{2\pi\i} = \\
&= \prod_{j=1}^N \int_{\mathbb{R}^2} \powe{-(v_\r^2 + v_\i^2) \lambda_j +  v_\r  2\Re (U^\dagger h)_j + \i v_\i 2\i\Im (U^\dagger h)_j } \frac{\mathrm{d}v_\r\mathrm{d}v_\i}{2\pi\i}.
\end{align}
At this point we may use 
\begin{align}
\intr{\powe{-ax^2 + bx}}{x} &= \sqrt{\frac{\pi}{a}} \powe{\frac{b^2}{4a}} \label{eqn:e to the minus a x squared plus b x}
\end{align}
to solve the integrals over $v_\r$ and $v_\i$ giving
\begin{align}
I &= \prod_{j=1}^N \frac{\pi}{\lambda_j} \frac{1}{2\pi\i} \powe{ \frac{1}{\lambda_j} \klb*{\Re (U^\dagger h)_j}^2 } \powe{ \frac{1}{\lambda_j} \klb*{-\Im (U^\dagger h)_j}^2 } = \frac{1}{(2\i)^N} \frac{1}{\det A} \exp \kla*{ \sum_{j=1}^N \frac{1}{\lambda_j} \abs*{(U^\dagger h)_j}^2  }.
\end{align}
The expression in the exponent is simply $\brc*{h}{U D^{-1}U^\dagger }{h} = \brc*{h}{A^{-1}}{h}$, which completes the proof (except for the normalization factor in front). \\
From this result one can also derive the important matrix elements
\begin{align}
\bra*{u_j^\ast u_k} &= \frac{\int_{\mathbb{C}^N} \exp\kla*{-\brc*{u}{A}{u} } u_j^\ast u_k \, \frac{\mathrm{d}u_1\mathrm{d}u_1^\ast}{2\pi\i} \dots \frac{\mathrm{d}u_N\mathrm{d}u_N^\ast}{2\pi\i}}{\int_{\mathbb{C}^N} \exp\kla*{-\brc*{u}{A}{u} } \, \frac{\mathrm{d}u_1\mathrm{d}u_1^\ast}{2\pi\i} \dots \frac{\mathrm{d}u_N\mathrm{d}u_N^\ast}{2\pi\i}} =  (A^{-1})_{jk}. \label{eqn:exp minus u ast A u uj ast uk over exp minus u ast A u}
\end{align}
This follows from realizing that the numerator may be expressed as the derivative of the denominator with respect to $-A_{jk}$.
We already solved that integral, and thus
\begin{align}
\bra*{u_j^\ast u_k} &= \frac{1}{I[h=0]} \frac{\partial }{\partial (-A_{jk})} I[h=0] = -\det A \frac{\partial }{\partial A_{jk}} \frac{1}{\det A} = \frac{1}{\det A} \frac{\partial }{\partial A_{jk}} \det A = \frac{\partial }{\partial A_{jk}} \log \det A.
\end{align}
Using $\log \det A = \Tr \log A$ and $\kla*{\frac{\partial }{\partial A_{jk}} \log A}_{nl} = \sum_{m=1}^N \kla*{\frac{\partial A_{nm}}{\partial A_{jk}}}_{nm} (A^{-1})_{ml} $ gives
\begin{align}
I_{jk} &= \Tr \kla*{ \sum_{m=1}^N  \delta_{jn} \delta_{mk} (A^{-1})_{ml} }_{nl} = \sum_{n=1}^N  \delta_{jn} (A^{-1})_{nk} = (A^{-1})_{jk}.
\end{align}
\subsection{Matsubara Summation}
We can evaluate the series of a function over even $\omega_n^\B = \frac{2\pi n}{\beta}$ or odd Matsubara frequencies $\omega_n^\F = \frac{(2n+1)\pi}{\beta}$ where $n\in\mathbb{Z}$ using the Residue theorem in reverse.
Denote the poles of $g(z)$ by $z_k$ and choose the contour $z=R\powe{\i t}$, $t\in [0,2\pi)$ for large $R$.
Since the Bose (Fermi) occupation function $f_{\B/\F}(z) = \frac{1}{\powe{\beta z} \mp 1}$ has poles at $\i\omega_n^{\B/\F}$ with residue $\lim_{z\rightarrow \i\omega_n^{\B/\F}} \frac{z - \i\omega_n^{\B/\F}}{\powe{\beta z } \mp 1} = \frac{1}{\beta} \powe{-\beta \i \omega_n^{\B/\F}} = \frac{\pm 1}{\beta}$ we have
\begin{align}
\oint f_{\B/\F}(z) g(z) \,\mathrm{d}z &= 2\pi\i \sum_{n\in\mathbb{Z}} \Res \klb*{f_{\B/\F} g}(\i\omega_n^{\B/\F}) + 2\pi\i \sum_{z_k} \Res \klb*{f_{\B/\F} g}(z_k) = \\
&= 2\pi\i \sum_{n\in\mathbb{Z}} \frac{\pm 1}{\beta} g(\i\omega_n^{\B/\F})  + 2\pi\i \sum_{z_k} f_{\B/\F}(z_k) \Res \klb*{g}(z_k).
\end{align}
Assuming the integral to vanish, which is justified for $|g(z)| \lesssim \frac{1}{|z|^{1+\epsilon}}$, then
\begin{align}
\sum_{n\in\mathbb{Z}} g(\i\omega_n^{\B/\F}) &= \mp \beta \sum_{z_k} f_{\B/\F}(z_k) \Res[g] (z_k). \label{eqn:sum g of i omega n BF is mp  beta sum fBF of zk times Res g of zk}
\end{align}
In general the expression on the right hand side will be significantly easier to evaluate.
However, one should be mindful of the convergence criteria to avoid mistakes introduced by this replacement.
In essence, we need to consider the inequality
\begin{align}
\abs*{\oint_{|z|=R} f_{\B/\F}(z) g(z) \,\mathrm{d}z} &\le \intx{0}{2\pi}{ \abs*{f_{\B/\F}\kla*{R\powe{\i t}} \frac{1}{|R\powe{\i t}|^{1+\epsilon}}} R }{t} = \frac{1}{R^\epsilon} \intx{0}{2\pi}{\frac{1}{ \abs*{\powe{\beta R \powe{\i t}} \mp 1} }}{t} \overset{R \rightarrow \infty}{\longrightarrow} 0.
\end{align}
\subsection{Correlations}
\subsubsection{Basic correlations}
Let $r=(x,u\tau)$ with some velocity $u$ and the imaginary time $\tau = \i t$.
Also let $\textbf{q} = \kla*{k, \frac{\omega_n}{u}}$ with the Matsubara frequencies $\omega_n$ and write $\textbf{q}r \equiv kx - \omega_n \tau$.
We define the Fourier transform as
\begin{align}
\phi(r) &= \sum_\textbf{q} \phi(\textbf{q}) \powe{\i r\textbf{q}},\\
\phi(\textbf{q}) &= \frac{1}{\beta L} \int \phi(r) \powe{-\i r\textbf{q}} \,\mathrm{d}\textbf{q}.
\end{align}
Then we can compute the correlation
\begin{align}
\bra*{\klb*{\phi(r_1) - \phi(r_2)}^2} &= \bra[\bigg]{ \klb[\bigg]{\sum_{\textbf{q} } \phi(\textbf{q}) \kla*{\powe{\i r_1\textbf{q}} - \powe{\i r_2\textbf{q}}}}^2   } = \\
&= \sum_{\textbf{q}_1,\textbf{q}_2} \bra*{\phi(\textbf{q}_1) \phi(\textbf{q}_2)}  \kla*{\powe{\i r_1\textbf{q}_1} - \powe{\i r_2\textbf{q}_1}} \kla*{\powe{\i r_1\textbf{q}_2} - \powe{\i r_2\textbf{q}_2}}.
\end{align}
The expectation values in momentum space can be represented as a path integral
\begin{align}
\bra*{\phi^\ast(\textbf{q}_1) \phi(\textbf{q}_2)} &= \frac{1}{Z} \int \phi^\ast(\textbf{q}_1) \phi(\textbf{q}_2) \powe{-\frac{S}{\hbar}} \,\mathcal{D}\phi(\textbf{q}_1) \mathcal{D}\phi(\textbf{q}_2)
\end{align}
where the action is
\begin{align}
-\frac{S}{\hbar} &= \intx{0}{\beta}{ \intx{ }{ }{ \klb*{\i\Pi \partial_\tau \phi - H(\phi,\Pi)}  }{x}}{\tau}.
\end{align}
For a typical quadratic Luttinger Hamiltonian and $\Pi = \frac{1}{\pi} \nabla \theta$ this gives
\begin{align}
-\frac{S}{\hbar} &= \intx{0}{\beta}{ \intx{ }{ }{ \klb*{\i\frac{1}{\pi} \nabla\theta \partial_\tau \phi - \frac{u}{2\pi} \kla*{K(\nabla\theta)^2 + \frac{1}{K} (\nabla\phi)^2} }  }{x}}{\tau}.
\end{align}
Using Fourier components this is equivalent to
\begin{align}
-\frac{S}{\hbar} &= \frac{1}{\pi}\sum_{\textbf{q},\textbf{q}'} \intx{0}{\beta}{\intx{ }{ }{ \klb*{ \omega_n \phi(\textbf{q}) \i k' \theta(\textbf{q}')  + \frac{u}{2} \kla*{K k k' \theta(\textbf{q}) \theta(\textbf{q}') + \frac{1}{K} kk' \phi(\textbf{q}) \phi(\textbf{q}')  }   } \powe{\i r\textbf{q}} \powe{\i r\textbf{q}'} }{x}}{\tau} = \\
&= -\frac{\beta L}{\pi} \sum_{\textbf{q}} \klb*{\omega_n \i k\phi(\textbf{q}) \theta(-\textbf{q}) + \frac{uk^2}{2} \kla*{K \theta(\textbf{q}) \theta(-\textbf{q}) + \frac{1}{K} \phi(\textbf{q}) \phi(-\textbf{q})} }.
\end{align}
Note that $\phi(-\textbf{q}) = \phi^\ast(\textbf{q})$ and similarly for $\theta$.
Then the action may also be written as
\begin{align}
\frac{S}{\hbar} &= \frac{\beta L}{2\pi} \sum_{\textbf{q}} \begin{pmatrix}
\theta_\textbf{q}^\ast, \phi_\textbf{q}^\ast
\end{pmatrix} \begin{pmatrix}
uKk^2 & \i\omega_n k \\ \i\omega_n k & \frac{u}{K}k^2
\end{pmatrix} \begin{pmatrix}
\theta_\textbf{q} \\ \phi_\textbf{q}
\end{pmatrix} =: \frac{\beta L}{2\pi} \sum_\textbf{q} \Gamma_\textbf{q}^\dagger M^{-1} \Gamma_\textbf{q}.
\end{align}
Since the quantity we want to average does not involve and $\theta$-terms it is possible to eliminate the corresponding integration.
We may write
\begin{align}
\frac{S}{\hbar} &= \frac{\beta L}{2\pi} \sum_\textbf{q} \klb*{\frac{uk^2}{K} \phi_\textbf{q} \phi_{-\textbf{q}} + uk^2 K \kla*{\theta_\textbf{q} + \frac{\i\omega_nk\phi_\textbf{q}}{uk^2 K}} \kla*{\theta_{-\textbf{q}} + \frac{\i\omega_nk\phi_{-\textbf{q}}}{uk^2 K}} - uk^2 K \frac{\i\omega_nk\phi_\textbf{q}}{uk^2 K} \frac{ \i\omega_nk \phi_{-\textbf{q}}}{uk^2 K} }.
\end{align}
Defining the shifted variable $\tilde{\theta}_\textbf{q} := \theta_\textbf{q} + \frac{\i\omega_n}{ukK} \phi_\textbf{q}$ simplifies this expression to
\begin{align}
\frac{S}{\hbar} &= \frac{\beta L}{2\pi} \sum_\textbf{q} \klb*{ \kla*{\frac{uk^2}{K} + \frac{\omega_n^2}{uK} } \phi_\textbf{q} \phi_{-\textbf{q}} + uk^2K \tilde{\theta}_\textbf{q} \tilde{\theta}_{-\textbf{q}} }.
\end{align}
Both the numerator as well as the denominator $Z$ will contain the same contribution from the $\tilde{\theta}$-integration which consequently drops out.
This effectively leaves
\begin{align}
\bra*{\phi^\ast(\textbf{q}_1) \phi(\textbf{q}_2)} &= \frac{\int \powe{-\frac{S_\phi}{\hbar}} \phi_{\textbf{q}_1}^\ast \phi_{\textbf{q}_2} \, D\phi_\textbf{q} D\phi_\textbf{q}^\ast }{\int \powe{- \frac{S_\phi}{\hbar}}  \, D\phi_\textbf{q} D\phi_\textbf{q}^\ast} = (A^{-1})_{\textbf{q}_1 \textbf{q}_2}
\end{align}
where the relevant matrix is implicitly defined via $\frac{S}{\hbar} = \sum_{\textbf{q}\textbf{q}'} \phi_\textbf{q}^\ast A_{\textbf{q} \textbf{q}'} \phi_{\textbf{q}'}$ and is explicitly given by
\begin{align}
A_{\textbf{q}\textbf{q}'} &= \delta_{\textbf{q}\textbf{q}'} \frac{\beta L}{2\pi uK} \kla*{u^2k^2 + (\omega_n^\B)^2}.
\end{align}
The expectation value is thus
\begin{align}
\bra*{\phi^\ast(\textbf{q}_1) \phi(\textbf{q}_2)} &= \delta_{\textbf{q}_1 \textbf{q}_2} \frac{2\pi uK}{\beta L} \frac{1}{u^2k^2 + (\omega_n^\B)^2}.
\end{align}
\TODO{additional factor of $2\pi$? $\hbar$ probably inconsistent?}\\
One may also make use of \autoref{eqn:exp minus u ast A u uj ast uk over exp minus u ast A u} more directly since $\bra*{\phi^\ast(\textbf{q}_1) \phi(\textbf{q}_2)}$ is the $\phi,\phi^\ast$ component of
\begin{align}
\kla*{\frac{\beta L}{2\pi} M^{-1}}^{-1} &= \frac{2\pi}{\beta L} \frac{1}{u^2 k^4 + (\omega_n^\B)^2 k^2 } \begin{pmatrix}
\frac{u}{K} k^2 & -\i\omega_n^\B k \\ -\i\omega_n^\B k & uKk^2
\end{pmatrix} = \frac{2\pi}{\beta L} \frac{1}{u^2k^2 + (\omega_n^\B)^2} \begin{pmatrix}
\frac{u}{K} & -\frac{\i\omega_n}{k} \\ -\frac{\i\omega_n^\B}{k} & uK
\end{pmatrix}. \label{eqn:M inverse k-correlations}
\end{align}
With this information we can continue our computation of the correlation function
\begin{align}
\bra*{\klb*{\phi(r) - \phi(0)}^2} &= \sum_{\textbf{q}} \frac{2\pi uK}{\beta L} \frac{1}{u^2k^2 + (\omega_n^\B)^2} \kla*{\powe{\i r\textbf{q}} - 1} \kla*{\powe{-\i r\textbf{q}} - 1} = \\
&= \frac{2\pi uK}{\beta L} \sum_k \sum_{n\in\mathbb{Z}} \frac{1}{u^2k^2 + (\omega_n^\B)^2} \kla*{2 - 2 \cos (kx - \omega_n^\B \tau)} = \\
&= \frac{uK}{\beta} \sum_{n\in\mathbb{Z}} \intx{}{}{\frac{2}{u^2k^2 + (\omega_n^\B)^2} \kla*{1 - \cos(kx - \omega_n^\B\tau)} }{k}.
\end{align}
%%\TODO{
%%Giamarchi claims divergence of integrals at large $k$, but why? 
%%These seem to converge (for simplicity assume $x\ge 0$, otherwise need additional care of residue signs)
%%\begin{align}
%%\dots &= \frac{2uK}{\beta} \sum_{\omega_n} \intr{\frac{1}{u^2k^2 + \omega_n^2}}{k} - \frac{2uK}{\beta} \sum_{\omega_n} \intr{\frac{\cos(kx - \omega_n\tau)}{u^2k^2 + \omega_n^2}}{k} = \\
%%&= \frac{2uK}{\beta} \sum_{\omega_n} \klb*{2\pi\i \frac{1}{u^2} \lim_{k \rightarrow \i \frac{\omega_n}{u}} \frac{k - \i \frac{\omega_n}{u}}{k^2 + \frac{\omega_n^2}{u^2}} - \Re\kla*{2\pi\i\frac{1}{u^2} \lim_{k \rightarrow \i \frac{\omega_n}{u}} \frac{k - \i \frac{\omega_n}{u}}{k^2 + \frac{\omega_n^2}{u^2}} \powe{\i (kx-\omega_n\tau)} } } = \\
%%&= \frac{2uK}{\beta} \sum_{\omega_n} \klb*{ \frac{2\pi\i}{u^2} \frac{1}{2\i \frac{\omega_n}{u}} - \Re \kla*{\frac{2\pi\i}{u^2} \frac{\powe{-\i\omega_n\tau} \powe{-\frac{\omega_n}{u}x} }{2\i \frac{\omega_n}{u}} } } = \\
%%&= \frac{2\pi K}{\beta} \sum_{\omega_n} \klb*{\frac{1}{\omega_n} - \frac{\cos (\omega_n\tau)}{\omega_n} \powe{-\frac{\omega_n}{u}x} }.
%%\end{align}
%%Okay, maybe the $\frac{1}{\omega_n}$-sum doesn't look that good... is this the issue (from a mathematical POV)?
%%}\\
To avoid potential problems with divergences let us introduce an additional factor of $\powe{-\alpha|k|}$ in the integrand.
Then the correlation function is
\begin{align}
\bra*{\klb*{\phi(r) - \phi(0)}^2} &= \frac{2K}{\beta u} \sum_{n\in\mathbb{Z}} \intr{\frac{\powe{-\alpha|k|}}{k^2 + \frac{(\omega_n^\B)^2}{u^2}}  \kla*{1 - \Re \powe{\i (kx -\omega_n^\B \tau)}} }{k} =: KF_1(r).
\end{align}
We can use the residue theorem to solve this integral.
Note that depending on the sign of $x$ the contour must be closed in the upper (lower) half of the complex plane when $x> 0$ ($x<0$).
Closing in the lower half plane also changes the global orientation of the contour which introduces an additional minus sign.
The relevant pole is thus $k_0 = \sgn (x) \i \frac{|\omega_n^\B|}{u}$ and we get
\begin{align}
F_1(r) &= \frac{2}{\beta u} \sum_{n\in\mathbb{Z}} 2\pi\i \sgn(x)  \lim_{k \rightarrow k_0} \frac{\powe{-\alpha|k|}}{k + k_0} \kla*{1 - \Re \powe{\i (kx-\omega_n\tau)}} = \\
&= \frac{2\pi}{\beta u} \sum_{n\in\mathbb{Z}} \sgn(x) \frac{u}{|\omega_n^\B|} \powe{-\alpha \frac{|\omega_n^\B|}{u}} \kla*{1 - \cos(\omega_n^\B\tau) \powe{-\frac{|\omega_n^\B|}{u}|x|} }.
\end{align}
%\TODO{
%The summand, as a function of $\i\omega_n$, has a single pole at the origin.
%Using \autoref{eqn:sum g of i omega n BF is mp  beta sum fBF of zk times Res g of zk} for the Bose case we find
%\begin{align}
%\bra*{\klb*{\phi(r) - \phi(0)}^2} &= 2\pi K \sgn (x) \lim_{z\rightarrow 0} z \frac{1}{|z|} \dots
%\end{align}
%Actually, there is no pole at zero, and thus the sum is just zero (which is obviously incorrect). What went wrong here?\\
%Try different approach: $\sum_{n\ge 1} g(\i\omega_n^B) = -\beta \sum_{z_l} \dots + \int gf_\B$.
%Since $g(z) = \frac{1}{|z|} \powe{-\alpha \frac{|z|}{u}} \kla*{1 - \cos(\i z\tau) \powe{-\frac{|x|}{u}|z|} } $
%\begin{align}
%\dots &= \intr{g(x) \frac{1}{\powe{\beta x} - 1}}{x} + \lim_{\epsilon \rightarrow 0} \intx{\pi}{0}{ g(\epsilon\powe{\i t}) \frac{1}{\powe{\beta \epsilon \powe{\i t}} - 1} \epsilon\i \powe{\i t}}{t} = \\
%&= \intr{\frac{g(x)}{\powe{\beta x} - 1}}{x} + \intx{\pi}{0}{ \frac{|x|}{u} \frac{\i\powe{\i t}}{\beta \powe{\i t}} }{t} = \intr{\frac{g(x)}{\powe{\beta x} - 1}}{x} -\i \frac{|x|}{\beta u}.
%\end{align}
%%#LUTTINGER
%%alpha=0.1
%%x=5.6
%%tau=-0.3
%%u=1.0
%%beta=1.0
%%N=1000
%%SIZE=100
%%omegan = 2*np.pi/beta * np.arange(-N, N+1)
%%omegan2 = np.copy(omegan[omegan > 0])
%%def corr(k, omega, x, tau, u, beta):
%%    return 2*u/beta * (1 - np.cos(k*x-omega*tau)) / (k**2*u**2 + omega**2) * np.exp(-alpha * np.abs(k))
%%
%%print(2*np.pi/beta * (np.sum(1/omegan2 * np.exp(-alpha/u * np.abs(omegan2)) * (1 - np.cos(omegan2 * tau) * np.exp(-np.abs(x) / u * np.abs(omegan2))) ) * 2 + x/u))
%%print(quad(lambda k: np.sum(corr(k, omegan, x, tau, u, beta)), -np.inf, np.inf))
%%giaint = lambda k: 2*np.exp(-alpha * k)/k * (2/(np.exp(beta*u*k) - 1) * (1 - np.cos(k*x) * np.cosh(tau*u*k)) + (1 - np.cos(k*x) * np.exp(-np.abs(tau*k*u))))
%%print(quad(giaint, 0, SIZE))
%}
%%%NEW APPROACH: Matsubara summation first.
%%%\begin{align}
%%%\bra*{\klb*{\phi(r) - \phi(0)}^2} &= \frac{2Ku}{\beta } \sum_{\omega_n} \intr{\frac{\powe{-\alpha|k|}}{u^2k^2 + \omega_n^2}  \kla*{1 - \Re \powe{\i (kx -\omega_n\tau)}} }{k} =: \intr{\sum_{n\in\mathbb{Z}} g(\i\omega_n^\B)}{k}.
%%%\end{align}
%%%Using \autoref{eqn:sum g of i omega n BF is mp  beta sum fBF of zk times Res g of zk} this gives\footnote{$f_\B(x) + f_\B(-x) \equiv -1$.}
%%%\begin{align}
%%%\bra*{\klb*{\phi(r) - \phi(0)}^2} &= \intr{ (-\beta)\sum_{z_l=\pm uk} \Res[g](z_l) f_\B(z_l) }{k} = \\
%%%&= -2Ku \intr{ \powe{-\alpha|k|}  \sum_{z_l=\pm uk} \frac{1}{-2z_l} f_\B(z_l) \kla*{1 - \powe{\i kx - z_l \tau} }  }{k} = \\
%%%&= K \intr{\frac{\powe{-\alpha |k|}}{k} \klb*{f_\B(uk) \kla*{1 - \powe {\i kx- ku\tau}} - f_\B(-uk) \kla*{1 - \powe{\i kx + ku\tau}} }  }{k} = \\
%%%&= 2K \intx{0}{\infty}{\frac{\powe{-\alpha|k|}}{k} \klb*{f_\B(uk) \kla*{1 - \cos(kx) \powe{-ku\tau}} - f_\B(-uk)\kla*{1 - \cos(kx) \powe{ku\tau}} } }{k} = \\
%%%&= 2K \intx{0}{\infty}{\frac{\powe{-\alpha|k|}}{k} \klb*{2f_\B(uk) \kla*{1 - \cos(kx) \cosh(ku\tau)} + \kla*{1 - \cos(kx) \powe{ku\tau}} } }{k}.
%%%\end{align}
%%%\TODO{This is correct except for $\tau \rightarrow -|\tau|$. Where can this replacement come from?}\\
%%%%%!!BREAK!!%%%%%
%NEW
%Note that $\cos(kx\pm \omega_n^\B \tau) = \cos(kx)\cos(\omega_n^\B\tau)$ under the integral, as the the antisymmetric term vanishes upon $k$-integration. 
%Then
%\begin{align}
%F_1(r) &= \frac{2u}{\beta} \sum_{n\in\mathbb{Z}} \intr{\powe{-\alpha|k|} \frac{1}{k^2u^2 + (\omega_n^\B)^2} \kla*{1 - \powe{\i kx} \cos(\omega_n^\B \tau) } }{k} = \\
%&= \frac{2}{\beta u} \sum_{n\in\mathbb{Z}} 2\pi\i \powe{-\alpha \frac{|\omega_n^\B|}{u}} \frac{1}{2\i \frac{|\omega_n^\B|}{u}} \kla*{1 - \powe{- x\frac{|\omega_n^\B|}{u}} \cos \kla*{\omega_n^\B\tau} } = \\
%&= \frac{2\pi}{\beta} \frac{x}{u} + \frac{2\pi}{\beta} 2\Re \sum_{n\ge 1} \powe{-\frac{\alpha}{u} \omega_n^\B} \frac{1}{\omega_n^\B} \kla*{1 - \powe{-\frac{x}{u} \omega_n^\B} \powe{\i \tau \omega_n^\B}}.
%\end{align}
Splitting of the $n=0$-term and using the symmetry we can reduce this expression to
\begin{align}
F_1(r) &= \frac{2\pi}{\beta} \frac{x}{u} + \sgn(x) \frac{2\pi}{\beta} 2\Re \sum_{n\ge 1} \powe{-\frac{\alpha}{u} \omega_n^\B} \frac{1}{\omega_n^\B} \kla*{1 - \powe{-\frac{|x|}{u} \omega_n^\B} \powe{\i \tau \omega_n^\B}}.
\end{align}
Consider the general series for some complex number $c\in\mathbb{C}$ with $\Re(c) > 0$
\begin{align}
S(c) &:= \sum_{n\ge 1} \frac{1}{n} \powe{-cn}.
\end{align}
The derivative with respect to $\Re(c)$ is a geometric series
\begin{align}
\frac{\partial S}{\partial \Re(c)} &= -\sum_{n\ge 1} \powe{-cn} = -\frac{\powe{-c}}{1 - \powe{-c}} = -\frac{1}{\powe{c} - 1}.
\end{align}
Using the limit $S(c\rightarrow \infty) \rightarrow 0$ gives
\begin{align}
S(c) &= \intx{\Re(c)}{\infty}{\frac{1}{\powe{\i \Im(c)} \powe{z} - 1}}{z} = \log \kla*{1 - \powe{-\i \Im(c)} \powe{-z}} \Big\vert_{\Re(c)}^\infty = -\log \kla*{1 - \powe{-c}}.
\end{align}
Also note that $2\Re \kla*{S(c)} = -\log \kla*{1 + \powe{-2\Re (c)} - 2\powe{-\Re(c)} \cos\kla*{\Im(c)} } $.
Our correlation function is equivalent to the difference of two such series since
\begin{align}
\sgn(x) F_1(r) &= \frac{2\pi |x|}{\beta u} + 2 \Re \klb*{S\kla*{\frac{2\pi\alpha}{\beta u}} - S\kla*{\frac{2\pi(\alpha + |x|)}{\beta u} + \frac{2\pi \tau}{\beta}\i} } = \\
&= \frac{2\pi |x|}{\beta u} - 2 \log \kla*{1 - \powe{-\frac{2\pi\alpha}{\beta u}}} +  \log \kla*{1 + \powe{-\frac{4\pi (\alpha + |x|)}{\beta u}} - 2\powe{-\frac{2\pi (\alpha + |x|)}{\beta u}} \cos \kla*{\frac{2\pi\tau}{\beta}} }.
\end{align}
This expression can be simplified a bit more by bringing all terms into one logarithm
\begin{align}
\powe{\sgn(x) F_1(r)} &= \frac{\powe{\frac{2\pi |x|}{\beta u}} + \powe{-\frac{2\pi (2\alpha  + |x|)}{\beta u}} - 2 \powe{-\frac{2\pi \alpha}{\beta u}} \cos \kla*{\frac{2\pi\tau}{\beta}} }{\kla*{1 - \powe{-\frac{2\pi\alpha}{\beta u}}}^2} = \\
&= \frac{2\cosh\kla*{\frac{2\pi (\alpha + |x|)}{\beta u}} - 2\cos \kla*{\frac{2\pi\tau}{\beta}} }{4\sinh^2\kla*{\frac{\pi\alpha}{\beta u}} } = \\
&= \frac{\sinh^2 \kla*{\frac{\pi (\alpha + |x|)}{\beta u}} + \sin^2\kla*{\frac{\pi\tau}{\beta}} }{\sinh^2\kla*{\frac{\pi\alpha}{\beta u}} }.
\end{align}
Thus the first special function is
\begin{align}
F_1(r) &= -2\sgn(x) \log \sinh \kla*{\frac{\pi\alpha}{\beta u}} +\sgn(x)  \log \klb*{\sinh^2 \kla*{\frac{\pi (\alpha + |x|)}{\beta u}} + \sin^2\kla*{\frac{\pi\tau}{\beta}}}.
\end{align}
\TODO{This is actually more accurate than the approximation given in Giamarchi (C.28)!
For reference, he gives $F_1(r) = \QnA{\frac{1}{2}} \log \klb*{\frac{\beta^2u^2}{\pi^2\alpha^2} \kla*{\sin^2 \kla*{\frac{\pi x}{\beta u}} + \sinh^2\kla*{\frac{\pi\tau}{\beta}}} } $..
However, the sign function is probably a mistake and I have an additional factor of 2...}\\
Following \autoref{eqn:M inverse k-correlations} we can compute the $\theta$-correlation function from the $\phi$-correlation function by the replacement $K\rightarrow \frac{1}{K}$, i.e. $\bra*{\klb*{\theta(r) - \theta(0)}^2} = \frac{1}{K}F_1(r)$.
The mixed correlations are
\begin{align}
\bra*{\phi(r_1)\theta(r_2)} &= \sum_{\textbf{q}_1, \textbf{q}_2} \bra*{\phi(\textbf{q}_1) \theta(\textbf{q}_2)} \powe{\i r_1\textbf{q}_1} \powe{\i r_2\textbf{q}_2} = \\
&= \frac{2\pi}{\beta L} \sum_{\textbf{q}} \frac{1}{k} \frac{-\i \omega_n^\B}{u^2k^2 + (\omega_n^\B)^2} \powe{\i (r_1-r_2)\textbf{q}} =: \frac{1}{2}F_2(r).
\end{align}
Again assuming $L\rightarrow \infty$ we have
\begin{align}
F_2(r) &= \frac{-2\i}{\beta u^2} \sum_{n\in\mathbb{Z}, n\neq 0} \intr{\frac{\powe{-\alpha |k|}}{k} \frac{\omega_n^\B}{k^2 + \frac{(\omega_n^\B)^2}{u^2}} \powe{\i (kx-\omega_n^\B \tau)} }{k} .
\end{align}
Similar to the previous case the relevant pole is $k_0 = \i \sgn(x) \frac{|\omega_n^\B|}{u}$.
An additional complexity is the need to avoid the origin.
This is done by a small semicircle parametrized by $k(t) = \epsilon \powe{\sgn(x) \i t}$ for $t\in [\pi, 0]$.
The corresponding integral gives
\begin{align}
I_\epsilon &= \frac{-2\i}{\beta u^2} \sum_{n\in\mathbb{Z}, n\neq 0} \sgn(x) \i \intx{\pi}{0}{ \powe{-\alpha \epsilon} \frac{\omega_n^\B}{\epsilon^2 \powe{ \sgn(x)2\i t} + \frac{(\omega_n^\B)^2}{u^2} } \powe{-\i \omega_n^\B\tau} \powe{\i x \epsilon \powe{\sgn(x) \i t}} }{t} \\
&\overset{\epsilon\rightarrow 0}{\longrightarrow} \frac{2}{\beta u^2} \sum_{n\in\mathbb{Z}, n\neq 0} \sgn(x) \omega_n^\B \powe{-\i\omega_n^\B\tau} \frac{u^2}{(\omega_n^\B)^2} \intx{\pi}{0}{ }{t} = \frac{2\pi}{\beta} \sgn(x) 2\i\Im\sum_{n\ge 1} \frac{\powe{\i\omega_n^\B\tau} }{\omega_n^\B}.
\end{align}
Using our formula gives
\begin{align}
I_{\epsilon=0} &= -\sgn(x) 2\i\Im \log \kla*{1 - \powe{\i \frac{2\pi\tau}{\beta}}}.
\end{align}
For $x<0$ the orientation of the contour changes which produces a minus sign in front of the residue giving
\begin{align}
F_2(r) &= -I_{\epsilon = 0} + 2\pi\i \sgn(x) \frac{-2\i}{\beta u^2} \sum_{n\in\mathbb{Z},n\neq 0} \omega_n^\B \frac{\powe{-\alpha|k_0|}}{2k_0^2} \powe{\i k_0x} \powe{-\i \omega_n^\B\tau} = \\
&= \sgn(x) 2\i\Im \log \kla*{1 - \powe{\i \frac{2\pi\tau}{\beta}}} + \sgn(x) \frac{2\pi}{\beta} 2\i\Im\sum_{n\ge 1} \frac{\powe{-\frac{\alpha}{u} \omega_n^\B}}{\omega_n^\B} \powe{\i\omega_n^\B\tau} \powe{-\frac{|x|}{u} \omega_n^\B} = \\
&= \sgn(x) 2\i\Im\klc*{\log \kla*{1 - \powe{\i\frac{2\pi\tau}{\beta}}} - \log \kla*{1 - \powe{-\frac{2\pi(\alpha + |x|)}{\beta u} + \i\frac{2\pi\tau}{\beta}}}  } = \\
&= \sgn(x) 2\i \Im \log \kla*{\frac{\powe{-\i\frac{\pi\tau}{\beta}} - \powe{\i \frac{\pi\tau}{\beta}} }{ \powe{-\i\frac{\pi\tau}{\beta}} - \powe{-\frac{2\pi(\alpha + |x|)}{\beta u}} \powe{\i \frac{\pi\tau}{\beta}} }}.
\end{align}
%Let us first compute 
%\begin{align}
%\sum_{n\ge 1} \frac{\powe{-\frac{2\pi(\alpha + |x|)}{\beta u} n}}{n} \powe{\i \frac{2\pi\tau}{\beta} n} &= -\log \kla*{1 - \powe{-\frac{2\pi (\alpha + |x|)}{\beta u} + \i \frac{2\pi\tau}{\beta}}}.
%\end{align}
%The imaginary part takes the form $\Im\kla*{-\log \kla*{1 - \powe{-c}}} = -\arg\kla*{1 - \powe{-c}}$.
%We can also write this as
%\begin{align}
%\arg\kla*{1 - \powe{-2(a+\i b)}} &= \arg \kla*{\powe{a} \powe{\i b} - \powe{-a} \powe{-\i b}} - b = \\
%&= \arg \kla*{2\sinh(a)\cos(b) + 2\cosh(a) \i\sin(b)} - b = \\
%&= \arg \kla*{\sinh(a) \cos(b) + \i \cosh(a) \sin(b)} - b = \\
%&= \arg \kla*{\tanh(a) + \i \tan(b)} - b.
%\end{align}
%In our case $a = \frac{\pi (\alpha + |x|)}{\beta u}$ and $b = -\frac{\pi\tau}{\beta}$
%\begin{align}
%F_2(r) &= -I_{\epsilon = 0} - 2\i \klc*{\arg \klb*{\tanh\kla*{\frac{\pi (\alpha + |x|)}{\beta u}} - \i \tan \kla*{\frac{\pi\tau}{\beta}} } + \frac{\pi\tau}{\beta} } = \\
%&= -2\i \arg \klb*{\tan \kla*{\frac{\pi\tau}{\beta}} + \i\tanh\kla*{\frac{\pi (\alpha + |x|)}{\beta u}}} - \frac{2\pi\i\tau}{\beta} + \pi\i - I_{\epsilon=0}.
%\end{align}
In general $\Im\log (z) = \arg(z)$ and 
\begin{align}
\arg \kla*{\frac{\powe{a}\powe{-\i b} - \powe{-a} \powe{\i b}}{\powe{-\i b} - \powe{\i b} }} &= \arg \kla*{\frac{2\sinh(a) \cos(b) - 2\cosh(a)  \i\sin(b)}{-2\i \sin(b)}} = \arg\kla*{\tan(b) + \i \tanh(a)}.
\end{align}
In our case $a = \frac{\pi (\alpha + |x|)}{\beta u}$ and $b = \frac{\pi\tau}{\beta}$
\begin{align}
F_2(r) &= -\sgn(x) 2\i \arg \klb*{\tan\kla*{\frac{\pi\tau}{\beta}} + \i \tanh \kla*{\frac{\pi (\alpha + |x|)}{\beta u}}}.
\end{align}
In this case the result converges for $\alpha \rightarrow 0$ and we get
\begin{align}
F_2(r) &= -2\i\arctan \klb*{\frac{\tanh \kla*{\frac{\pi x}{\beta u}} }{\tan \kla*{\frac{\pi\tau}{\beta }} }} .
\end{align}
\TODO{Up to the same factor of 2 this is the so called ``naive'' result given in Giamarchi (C.28).
However, numerically this seems to be identical for $\alpha \rightarrow 0$ so I do not get the subtlety here.
Of course the result is only well-defined modulo $2\pi$ and as such replacing $\arg$ with $\arctan$ can in principle differ.}\\
REMARK: Subtlety in $F_2$.\\
Explained in Giamarchi (C.40)--(C.45).
\subsubsection{General correlations}
Let $\phi_j$ for $j=1,\dots,N$ be bosonic operators.
Then we can derive a formula for the general expectation value
\begin{align}
\bra*{\prod_j \powe{\phi_j}} &= \bra*{\powe{\sum_j \phi_j}} \powe{\frac{1}{2} \sum_{j<k} \klb*{\phi_j, \phi_k}}.
\end{align}
For this we make use of the relation
\begin{align}
\bra*{\powe{\sum_j \phi_j}} &= \powe{\frac{1}{2} \bra*{\kla*{\sum_j \phi_j}^2}}.
\end{align}
Then we find
\begin{align}
\bra*{\prod_j \powe{\phi_j}} &= \powe{\frac{1}{2} \sum_j \bra*{\phi_j^2} + \frac{1}{2} \sum_{j<k} \bra*{\phi_j\phi_k} + \frac{1}{2} \sum_{j>k} \bra*{\phi_j\phi_k}} \powe{\frac{1}{2} \sum_{j<k} \bra*{\klb*{\phi_j, \phi_k}}} = \\
&= \powe{\frac{1}{2} \sum_j \bra*{\phi_j^2} + \sum_{j<k} \bra*{\phi_j \phi_k}}.
\end{align}
The following derivation for the general correlation function can be simplified by making use of this formula, translational invariance and $KF_1(r) = \bra*{\klb*{\phi(r) - \phi(0)}^2}$ and similar relations for the other expectation values.\\
But this is not required, one can also arrive at the correct result in a more brute force manner.
This will be shown next.\\
Now consider the general correlation function
\begin{align}
\bra*{\exp\kla[\Big]{\sum_j \i\kla*{A_j \phi(r_j) + B_j\theta(r_j)} } } &= \bra*{\exp\kla[\Big]{\i \sum_j \sum_\textbf{q} \kla*{A_j\phi(\textbf{q}) + B_j \theta(\textbf{q})} \powe{\i r_j\textbf{q}} }}.
\end{align}
Defining $A(\textbf{q}) = \sum_j A_j \powe{-\i r_j\textbf{q}}$ and similarly $B(\textbf{q})$ this gives
\begin{align}
\bra*{\exp \kla[\Big]{\sum_\textbf{q} \kla*{A(\textbf{q}) \phi(-\textbf{q}) + B(\textbf{q}) \theta(-\textbf{q})}} }  .
\end{align}
Inserting $\Ket{h} = -\i\kla[\big]{B(\textbf{q}), A(\textbf{q})}^T$, $\Ket{u} = \kla[\big]{\theta(\textbf{q}), \phi(\textbf{q})}^T$ and $A= \frac{\beta L}{2\pi} M^{-1}$ into \autoref{eqn:exp minus u ast A u plus h ast u plus u ast h} yields
\begin{align}
\kla*{A^{-1}} &= \exp\kla[\Big]{-\sum_\textbf{q} \brc*{h}{A^{-1}}{h} } = \\
&= \exp\kla*{-\frac{2\pi}{\beta L} \sum_{\textbf{q}} \begin{pmatrix}
B(-\textbf{q}),& A(-\textbf{q})
\end{pmatrix} \frac{1}{k^2u^2 + (\omega_n^\B)^2} \frac{1}{k} \begin{pmatrix}
\frac{u}{K}k & -\i\omega_n^\B \\ -\i\omega_n^\B & uKk
\end{pmatrix} \begin{pmatrix}
B(\textbf{q}) \\ A(\textbf{q})
\end{pmatrix} }.
\end{align}
The $A^2$-term is
\begin{align}
-\frac{2\pi}{\beta L} \sum_{\textbf{q}} \frac{uK}{k^2u^2 + (\omega_n^\B)^2} A(\textbf{q})A^\ast(\textbf{q}) = -\frac{2\pi}{\beta L} \sum_{\textbf{q}} \sum_{jj'} \frac{uK}{k^2u^2 + (\omega_n^\B)^2} A_jA_{j'} \cos\kla*{\textbf{q}[r_j - r_{j'}]} .
\end{align}
Replacing $\cos \rightarrow \cos -1 + 1$ recovers $F_1$ and allows us to write
\begin{align}
\frac{1}{2} \sum_{jj'} A_jA_{j'} KF_1(r_j - r_{j'}) - \kla[\Big]{\sum_j A_j}^2 \sum_\textbf{q} \frac{uK}{k^2u^2 + (\omega_n^\B)^2}.
\end{align}
Since the correlation function is the exponential of this expression and because the second term diverges to $+\infty$ we only get a non-zero result if $\sum_jA_j = 0$.
Likewise the $B^2$-terms require $\sum_jB_j = 0$ since the corresponding term can be found by the replacement $K\rightarrow\frac{1}{K}$.
The $AB$-terms give
\begin{align}
\frac{2\pi\i}{\beta L} \sum_{\textbf{q}} \klb*{B(-\textbf{q})A(\textbf{q}) + A(-\textbf{q})B(\textbf{q})} \frac{1}{k} \frac{\omega_n^\B}{k^2u^2 + (\omega_n^\B)^2} &= -\sum_{jj'} \klb*{A_jB_{j'} + B_j A_{j'}} \frac{1}{2} F_2(r_j - r_{j'})
\end{align}
And thus the full correlation function becomes
\begin{align}
\bra*{\powe{\sum_j \i\kla*{A_j \phi(r_j) + B_j\theta(r_j)} } } &= \powe{\frac{1}{2} \sum_{jj'} \klc*{\klb*{KA_jA_{j'} + \frac{1}{K}B_jB_{j'}} F_1(r_j - r_{j'}) - \klb*{A_jB_{j'} + B_jA_{j'} } F_2(r_j - r_{j'})}  }. \label{eqn:general correlation function}
\end{align}
\TODO{Again correct result up to factors of 2, \QnA{most likely from a misunderstanding in \autoref{eqn:exp minus u ast A u uj ast uk over exp minus u ast A u}}.
The book has an additional $\frac{1}{2}$ in the exponent, which is incompatible with my result.}
\subsubsection{Chrial fields}
The chiral fields are $\phi_\eta = K\theta - \eta \phi$ where $\eta = R,L$ or $\eta=\pm$ respectively.
The action can be expressed in these fields using the inverse transformation $\phi = \frac{\phi_L - \phi_R}{2}$ and $\theta = \frac{\phi_L + \phi_R}{2K}$
\begin{align}
\frac{S}{\hbar} &= \frac{\beta L}{2\pi} \sum_\textbf{q} \begin{pmatrix}
\theta_\textbf{q}^\ast, & \phi_\textbf{q}^\ast
\end{pmatrix} k\begin{pmatrix}
uKk & \i\omega_n^\B \\ \i\omega_n^\B & \frac{u}{K}k
\end{pmatrix} \begin{pmatrix}
\theta_\textbf{q} \\ \phi_\textbf{q}
\end{pmatrix} = \\
&= \frac{\beta L}{2\pi} \sum_\textbf{q} \begin{pmatrix}
\phi_{R,\textbf{q}}^\ast, & \phi_{L,\textbf{q}}^\ast
\end{pmatrix} \frac{1}{2K} \begin{pmatrix}
1 & -K \\ 1 & K
\end{pmatrix} k\begin{pmatrix}
uKk & \i\omega_n^\B \\ \i\omega_n^\B & \frac{u}{K}k
\end{pmatrix} \frac{1}{2K} \begin{pmatrix}
1 & 1 \\ -K & K
\end{pmatrix} \begin{pmatrix}
\phi_{R,\textbf{q}} \\ \phi_{L,\textbf{q}}
\end{pmatrix} = \\
&= \frac{\beta L}{2\pi} \sum_\textbf{q} \begin{pmatrix}
\phi_{R,\textbf{q}}^\ast, & \phi_{L,\textbf{q}}^\ast
\end{pmatrix} \frac{k}{2K} \begin{pmatrix}
uk - \i\omega_n^\B & 0 \\ 0 & uk + \i\omega_n^\B
\end{pmatrix}  \begin{pmatrix}
\phi_{R,\textbf{q}} \\ \phi_{L,\textbf{q}}
\end{pmatrix}.
\end{align}
\subsection{Analytical continuation}
In Fourier space the analytical continuation of a time-ordered Green's function is simply done by $G(\i\omega_n) \rightarrow G(\omega + \i\delta)$.\\
It is simpler to first obtain the ordered correlation function in \emph{real} time 
\begin{align}
\chi_{BA}^\t(t) &= - \Theta(t) \bra*{B(t)A(0)} \mp \Theta(-t) \bra*{A(0)B(t)}
\end{align}
where the upper sign corresponds to bosons and the lower sign corresponds to fermions.
$A(t) = \powe{\i Ht}A\powe{-\i Ht}$ is the usual Heisenberg evolution of an operator $A$. 
Note that
\begin{align}
\Theta(t) \chi_{BA}^\t(t) &= - \bra*{B(t)A(0)},\\
\Theta(t) \kla[\big]{\chi_{A^\dagger B^\dagger}^\t (-t)}^\ast &= \mp \bra*{B^\dagger(0) A^\dagger (-t)}^\ast = \mp \bra*{A(0)B(t)}.
\end{align}
The retarded correlation function is given by
\begin{align}
\chi_{BA}^\r(t) &= -\i \Theta(t) \bra*{\klb*{B(t)A(0)}_\mp} = \i\Theta(t) \klb*{\chi_{BA}^\t(t) - \kla[\big]{\chi_{A^\dagger B^\dagger}^\t (-t)}^\ast}.
\end{align}
One can obtain the time-ordered correlation functions in real time by performing the Wick rotation 
\begin{align}
\tau &= \i t + \epsilon \sgn(t) \label{eqn:Wick rotation}
\end{align}
in the time-ordered functions for imaginary time where $\epsilon=0^+$.
\subsubsection{Example: correlation of exponential of $\phi$}
Consider the correlation function
\begin{align}
\chi(x,\tau) &= -\bra*{\mathcal{T}_\tau \powe{\i 2\phi(x,\tau)} \powe{-\i 2\phi(0,0)} }.
\end{align}
This corresponds to \autoref{eqn:general correlation function} with $B_j\equiv 0$ and $A_{1,2} = \pm 2$ giving
\begin{align}
\chi(x,\tau) &= -\powe{\frac{1}{4} \sum_{jj'} \klc*{KA_jA_{j'} F_1(r_j - r_{j'}) )}  } = -\powe{-K F_1(r)} = \\
&= -\frac{\sinh^{2K} \kla*{\frac{\pi\alpha}{\beta u}} }{\klb*{\sinh^2\kla*{\frac{\pi (\alpha + |x|)}{\beta u}} + \sin^2\kla*{\frac{\pi\tau}{\beta}} }^K}.
\end{align}
For small $\alpha \ll |x|$ expanding to lowest order yields\footnote{Note $\sinh(a+\i b) \sinh(a-\i b) = \frac{1}{4} \kla*{\powe{a}\powe{\i b} - \powe{-a} \powe{-\i b}} \kla*{\powe{a} \powe{-\i b} - \powe{-a} \powe{\i b}} = \frac{\cosh (2a)}{2}  - \frac{\cos (2b)}{2} = \sinh^2 (a) + \sin^2(b)$.}
\begin{align}
\chi(x,\tau) &\simeq - \frac{\kla*{\frac{\pi\alpha}{\beta u}}^2}{\klb*{ \sinh^2\kla*{\frac{\pi x}{\beta u}} + \sin^2 \kla*{\frac{\pi\tau}{\beta}} }^K} = -\frac{\kla*{\frac{\pi\alpha}{\beta u}}^2}{\klb*{ \sinh\kla*{\frac{\pi}{\beta u} \klb*{x + \i\tau u} } \sinh\kla*{\frac{\pi}{\beta u} \klb*{x - \i\tau u} } }^K}.
\end{align}
The real time correlation function is thus
\begin{align}
\chi^\t(x,t) &= -\frac{\kla*{\frac{\pi\alpha}{\beta u}}^2}{\klb*{ \sinh\kla*{\frac{\pi}{\beta u} \klb*{x - tu + \i\epsilon u\sgn (t)} } \sinh\kla*{\frac{\pi}{\beta u} \klb*{x + tu - \i\epsilon u\sgn (t)} } }^K} \\
&=: -A^{-K} \kla*{\frac{\pi\alpha}{\beta u}}^{2K}.
\end{align}
\TODO{In this case $\chi_{A^\dagger B^\dagger}(\tau) = \chi_{AB}(\tau)$ (why is that? Is there a complex conjugation missing?)}
Note that the $\Theta(t)$-factor ensures $\sgn(t) = 1$ and we can therefore rewrite the complex part as $A^{-K} = \powe{-K\log (A)}$ where\footnote{Note $\sinh(a + \i\epsilon) \sinh(b - \i\epsilon) = \klb*{\sinh(a) + \i\epsilon \cosh (a)} \klb*{\sinh(b) - \i\epsilon \cosh(b)} + \mathcal{O}(\epsilon^2)$.}
\begin{align}
A &= \sinh\kla*{\frac{\pi}{\beta u} [x+tu]} \sinh\kla*{\frac{\pi}{\beta u}[x-tu]} + \i\epsilon \frac{\pi}{\beta} \sinh \kla*{\frac{2\pi t}{\beta}}. 
\end{align}
When $\epsilon \rightarrow 0$ we require the first term to be negative for $\log (A)$ to have an imaginary part.
This is the case when $(x+tu)(x-tu) < 0$, or equivalently, when $|x| < tu$.
In particular, since $\epsilon = 0^+$, the imaginary part of the logarithm will be positive, since $A$ lies in the upper half plane and we make a branch cut on the negative real axis.
Thus we can write
\begin{align}
\powe{-K\log (A)} \equiv \powe{- K\log (|A|) } \powe{-\i K\pi \Theta(tu - |x|)}
\end{align}
and the retarded correlation function becomes
\begin{align}
\chi^\r(t) &= -\Theta(t) 2\Im\chi^\t (x,t) = \\
&= -\Theta(t) \Theta (tu - |x|) 2\sin(K\pi) \kla*{\frac{\pi\alpha}{\beta u}}^{2K} \abs*{\sinh\kla*{\frac{\pi}{\beta u} [x+tu]} \sinh\kla*{\frac{\pi}{\beta u}[x-tu]}}^{-K}.
\end{align}
Define $\xi_\pm = tu \pm x$.
Performing the Fourier transform then gives
\begin{align}
\chi^\r(k,\omega) &= \intrndpi{2}{\powe{-\i (kx - \omega t)} \chi^\r(x,t)}{(x,t)} = \\
&= -\intx{0}{\infty}{\intx{-tu}{tu}{ 2\sin(K\pi) \kla*{\frac{\pi\alpha}{\beta u}}^{2K} \abs*{\sinh\kla*{\frac{\pi\xi_+}{\beta u}} \sinh\kla*{\frac{\pi\xi_-}{\beta u}}}^{-K} \powe{-\i (kx - \omega t)} }{x}}{t}.
\end{align}
We can transform to $\xi_\pm$ using $ \begin{pmatrix}
x \\ t
\end{pmatrix} = \frac{1}{2u}\begin{pmatrix}
u & -u \\ 1 & 1
\end{pmatrix} \begin{pmatrix}
\xi_+ \\ \xi_-
\end{pmatrix}$.
The determinant of this transformation is $\frac{1}{2u}$ and the region of integration is determined by $t> 0$ and $|x|< tu$.
Rewritten in terms of the new coordinates these are $\xi_++\xi_- > 0$ and $|\xi_+ - \xi_-| < \xi_+ + \xi_- $.
Neither condition places an upper bound on the variables, but they can not be negative.
We can express this change of variables as
\begin{align}
\intx{0}{\infty}{\intx{-tu}{tu}{ }{x}}{t} &= \frac{1}{2u} \intx{0}{\infty}{ }{\xi_+} \intx{0}{\infty}{ }{\xi_-}
\end{align}
and thus the correlation function becomes
\begin{align}
\chi^\r(k,\omega) &= -\frac{1}{u} \kla*{\frac{\pi\alpha}{\beta u}}^{2K} \sin(K\pi)  I\kla*{\frac{\omega - ku}{2u}, \frac{\pi}{\beta u}} I\kla*{\frac{\omega + ku}{2u}, \frac{\pi}{\beta u}}
\end{align}
where we defined the general integral
\begin{align}
I(a, b) &= \intx{0}{\infty}{ \powe{\i ax } \frac{1}{\sinh^K\kla*{bx} } }{x} = 2^K \intx{0}{\infty}{\powe{\i ax} \frac{\powe{- Kbx}}{\kla*{1 - \powe{- 2bx}}^K} }{x}.
\end{align}
Defining $y = 1 - \powe{-2bx}$ gives
\begin{align}
I(a,b) &= 2^K \intx{0}{1}{y^{-K} \powe{ (\i a - Kb) \frac{-1}{2b}\log (1-y)} \frac{1}{2b(1-y)} }{y} = \\
&= 2^K \frac{1}{2b} \intx{0}{1}{y^{-K} (1-y)^{\frac{K}{2} - 1 - \i\frac{a}{2b}} }{y} = 2^K \frac{1}{2b} B\kla*{1-K, \frac{K}{2} - \i \frac{a}{2b}}
\end{align}
where $B(x,y) = \frac{\Gamma(x)\Gamma(y)}{\Gamma(x+y)}$ is the Beta function.
The correlation function then becomes
\begin{align}
\chi^\r(k,\omega) &= -\frac{\sin(K\pi)}{u} \kla*{\frac{\pi\alpha}{\beta u}}^{2K}  4^K \frac{\beta^2 u^2}{4\pi^2} B\kla*{1 - K, \frac{K}{2} - \i \frac{\beta u}{2\pi} \frac{\omega + ku}{2u}} B\kla*{1 - K, \frac{K}{2} - \i \frac{\beta u}{2\pi} \frac{\omega - ku}{2u}} = \\
&= -\frac{\alpha^2\sin(K\pi)}{u} \kla*{\frac{2\pi\alpha}{\beta u}}^{2K-2} B\kla*{1 - K, \frac{K}{2} - \i \frac{\beta (\omega + ku)}{4\pi}} B\kla*{1 - K, \frac{K}{2} - \i \frac{\beta (\omega - ku)}{4\pi}}.
\end{align}
\subsection{Bosonization dictionary}
\section{Articles}
\subsection{245414 (2016): Dynamic response functions and helical gaps in Rashba nanowires}
\subsubsection{Bosonization}
The kinetic energy term is
\begin{align}
H_0 &= \sum_{p,s} \omega(p,s) \psi_{ps}^\dagger \psi_{ps}
\end{align}
with the dispersion $\omega(p,s) = \frac{p^2}{2m} - g_\text{R} s p$.
Let us first consider a simple term $\omega = p$ and drop the spin index.
Then we find
\begin{align}
H_0 &= \sum_{p} \frac{1}{L} \intr{\intr{\psi^\dagger(x') \psi(x) p\powe{\i p(x'-x)} }{x} }{ x'} = \\
&= \intr{\intr{\psi^\dagger(x') \psi(x) \i\nabla_x \delta(x-x') }{x} }{x'} = \\
&= \intr{\psi^\dagger(x) (-\i\nabla_x) \psi(x)}{x}.
\end{align}
Using the bosonization identities\footnote{$(\ell s)$ refers to the product as a single index}
\begin{align}
\psi_{\ell s}(x) &= \frac{\eta_{\ell s}}{\sqrt{2\pi\alpha}} \powe{\i\ell k_F x} \powe{-\i\ell \phi_{(\ell s)}(x) + \i\theta_{(\ell s)}(x)}
\end{align}
and $\psi_s = \psi_{\R s} + \psi_{\L s}$ we find (dropping the Klein-factors and all oscillatory or constant terms)
\begin{align}
\psi_s^\dagger (-\i\nabla_x) \psi_s &= \kla*{\powe{-\i k_Fx}\powe{\i\phi_s - \i\theta_s} + \powe{\i k_Fx} \powe{-\i\phi_{-s} - \i\theta_{-s}} } \frac{1}{2\pi\alpha} (-\i\nabla_x) \kla*{\powe{\i k_Fx}\powe{-\i\phi_s + \i\theta_s} + \powe{-\i k_Fx} \powe{\i\phi_{-s} + \i\theta_{-s}} } = \\
&= \frac{1}{2\pi\alpha} \kla*{ \nabla_x\klb*{-\phi_s +\theta_s } + \nabla_x \klb*{\phi_{-s} + \theta_{-s}} }.
\end{align}
Similarly, the second order derivatives yield
\begin{align}
\psi_s^\dagger (-\i\nabla_x)^2 \psi_x &= \frac{1}{2\pi\alpha} \kla*{\klb*{-\nabla_x\phi_s + \nabla_x \theta_s}^2 + \klb*{\nabla_x\phi_{-s} + \nabla_x\theta_{-s}}^2 + \nabla_x^2 \klb*{-\phi_s +\theta_s } + \nabla_x^2 \klb*{\phi_{-s} + \theta_{-s}}}.
\end{align}
Summing over $s$ results in the identities
\begin{align}
\sum_s \psi_s^\dagger (-\i\nabla_x)^2 \psi_s &= \frac{1}{2\pi\alpha} \kla*{2 (\nabla_x\phi_+)^2 + 2 (\nabla_x\theta_+)^2 + 2(\nabla_x\phi_-)^2 + 2(\nabla_x\theta_-)^2 + 2\nabla_x^2 \theta_+ + 2\nabla_x^2 \theta_-},\\
\sum_s s\psi_s^\dagger (-\i\nabla_x) \psi_s &= \frac{1}{2\pi\alpha} \kla*{2\nabla_x \phi_- - 2\nabla_x \phi_+}.
\end{align}
Upon integration over $x$ the total derivative terms drop out and we are left with
\begin{align}
H_0 &= \frac{u_0}{2\pi} \sum_s \intr{ \klb*{ (\nabla_x\theta_s)^2 + (\nabla_x\phi_s)^2 }  }{x}
\end{align}
where $u_0 = \frac{1}{\alpha m}$.
In the charge-spin basis $\phi_{\rho,\sigma} = \frac{\phi_+ \pm \phi_-}{\sqrt{2}}$ and $\theta_{\rho,\sigma} = \frac{\theta_+ \pm \theta_-}{\sqrt{2}}$ this becomes
\begin{align}
H_0 &= \frac{u_0}{2\pi} \sum_{j=\rho,\sigma} \intr{ \klb*{\kla*{\nabla_x\theta_j}^2 + \kla*{\nabla_x\phi_j}^2} }{x}.
\end{align}
\TODO{Including the interaction renormalizes the parameters, which introduces $U_\rho, K_\rho, u_\sigma, K_\sigma$.}\\
Additionally, there exist spin-exchange and spin-flip terms.
The former is given by
\begin{align}
V_\text{sx} &\propto \frac{1}{L}\sum_s \sum_{pp'q} \psi_{\L, s,p+q}^\dagger \psi_{\R, -s,p'-q}^\dagger \psi_{\R, s,p'} \psi_{\L, -s,p} = \\
&= \frac{1}{L^3} \sum_s \intrnd{4}{ \sum_{pp'q} \psi_{\L, s}^\dagger (x_1) \powe{\i x_1(p+q)} \psi_{\R, -s}^\dagger(x_2) \powe{\i x_2(p'-q)} \psi_{\R, s}(x_3) \powe{-\i x_3p'} \psi_{\L, -s}(x_4) \powe{-\i x_4 p} }{x} = \\
&= \sum_s \intr{\psi_{\L, s}^\dagger (x) \psi_{\R, -s}^\dagger(x) \psi_{\R, s}(x) \psi_{\L, -s}(x)}{x} \propto \\
&\propto \sum_s \intr{ \powe{-\i\phi_{-s} - \i\theta_{-s}} \powe{\i\phi_{-s} - \i\theta_{-s}} \powe{-\i\phi_s + \i\theta_s} \powe{\i\phi_s + \i\theta_s} }{x} = \\
&= \sum_s \intr{ \powe{2\i (\theta_s - \theta_{-s})}}{x} \propto g_\text{sx} \intr{\cos \kla[\big]{2\sqrt{2}\theta_\sigma} }{x}.
\end{align}
Similarly the spin-flip term becomes
\begin{align}
V_\text{sf} &\propto \frac{1}{L} \sum_s \sum_{pp'q} \psi_{\L, +,p+q}^\dagger \psi_{\L, +,p'-q}^\dagger \psi_{\R, -,p'} \psi_{\L, -,p} + (\L+ \leftrightarrow \R-) = \\
&= \intr{\psi_{\L, +}^\dagger (x) \psi_{\L, +}^\dagger(x) \psi_{\R, -}(x) \psi_{\R, -}(x)}{x} + (\L+ \leftrightarrow \R-) \propto \\
&\propto \sum_s \intr{ \klb*{\powe{-2\i\phi_{-} - 2\i\theta_{-}} \powe{-2\i\phi_- + 2\i\theta_- } + \text{h.c.}} }{x} = \\
&= \intr{ \powe{-4\i \phi_-}}{x} \propto g_\text{sf} \intr{\cos \kla[\big]{2\sqrt{2}[\phi_\rho - \phi_\sigma]} }{x}.
\end{align}
\TODO{(12) and (19) in article contains momentum prefactors $(2p'-q)(2p+q)$, why can those be assumed constant and absorbed in the coupling?
Or do these terms generate additional derivatives that can somehow be absorbed in the Luttinger parameters?
Also, global minus sign missing (maybe from Klein-factors)?}



%\begin{thebibliography}{9}
%\bibitem{example} description, \url{https://www.exmapleurl}, day.\,month~2021.
%\end{thebibliography}

\end{document}